For each integer $h\ge1$, let $T_h=\mathbb N^{\le h}$ be the
rooted tree of words of length at most $h$, ordered by extension of
words. Its root is the empty word. Keep separate copies of these trees
in the forest
\[
 \mathcal F=\coprod_{h\ge1}(\{h\}\times T_h).
\]
Write $a\preceq b$ if $a$ is an ancestor of $b$ in the same component.
A \emph{segment} is a nonempty interval
$[a,b]=\{v:a\preceq v\preceq b\}$ in an ancestral chain.
For a finitely supported real function $u$ on $\mathcal F$, set
$u(S)=\sum_{v\in S}u(v)$ and define
\begin{equation}\label{eq:segment-norm}
 \rho(u)=\sup\left\{\left(\sum_{i=1}^m u(S_i)^2\right)^{1/2}:
  m\ge1,\ S_1,\ldots,S_m\text{ pairwise disjoint segments}\right\}.
\end{equation}
The Euclidean triangle inequality and singleton segments show that
$\rho$ is a norm, with $\|u\|_2\le\rho(u)\le\|u\|_1$.
Let $J$ be its completion and let $X=J^*$, with
the dual norm denoted by $\|\cdot\|$. We write $e_v\in J$ for a
coordinate unit vector. Singleton tests give $|u(v)|\le\rho(u)$,
so coordinate evaluation extends to an element $e_v^*\in X$.
The construction in \eqref{eq:segment-norm} uses James's
square-sum norm on disjoint segments \cite[p.~739]{James1974}; the
branching and component heights here differ from the classical binary
tree.
